% --------------------------------- Preamble -----------------------------------
\documentclass[11pt, letterpaper, twocolumn]{article}

\usepackage[T1]{fontenc} % use 8-bit fonts
\usepackage[margin=0.75in, columnsep=0.2in, % 3.4 in columns
    headsep=0.28in, footskip=0.42in]{geometry} % paper sizes
\usepackage{amsmath, amsfonts, amssymb} % improved math
\usepackage{bm} % bold math (must be after amsmath)
\usepackage{upquote} % upright single quotes in verbatim
\usepackage{pgfplots} % includes tikz, xcolor, graphicx
\usepackage[colorlinks=true, allcolors=gray]{hyperref}
\usepackage{float}

\raggedbottom % do not force stretching text to fill the page
\setlength{\emergencystretch}{\hsize} % avoid overfull lines
\overfullrule=1mm % black band to show overfull lines
\pdfsuppresswarningpagegroup=1 % remove pdf image warnings
\frenchspacing % consistent spacing between words
\newcommand{\ul}[1] % vectors
    {{}\mkern1mu\underline{\mkern-1mu#1\mkern-1mu}\mkern1mu}
\DeclareSymbolFont{euler}{U}{eur}{m}{n} % symbol font
\DeclareMathSymbol{\PI}{\mathord}{euler}{25} % up pi
\pgfplotsset{compat=newest} % use newest pgfplot version
\newcommand{\EE}[2]{\ensuremath{{{#1}\!\times\!10^{#2}}}}
% ------------------------------------------------------------------------------

\title{RLC Circuit Simulation and State-Space Modeling}
\author{}
\date{}

\begin{document}
\maketitle

\section*{Description}
This project demonstrates the simulation of a continuous-time dynamical system
using Python. An RLC circuit is modeled using state-space equations and 
numerically integrated using a forward Euler method. This type of simulation
framework forms the basis for more advanced bayesian estimation techniques
such as Kalman Filtering.

We will see how we simulate the following RLC circuit: \\
\includegraphics[width=\linewidth]{../figures/fig_circuit.pdf}

\subsection*{System Model}
The system is modeled as an RLC circuit with the inductor current and capacitor
voltage chosen as the state variables:
\[
    \dot{x} = \begin{bmatrix}
        \dot{i}_L \\
        \dot{v}_C
    \end{bmatrix}  
\]
Using standard circuit relationship, the dynamics can be written as
\[
\dot{i}_L = \frac{1}{L}(v_s-v_C), \quad \dot{v_C} = \frac{1}{C}(i_L-\frac{v_C}{R})
\]
These equations describe the time evolution of the syste in terms of the 
current state and input voltage. For convenience, the system is expressed
in compact form as
\[
\dot{x} = f(x,u)
\]
where $u = v_s$ represents the input to the system.

\subsection*{Numerical Simulation}
The continuous-time dynamics are simulated using a forward Euler discretization
with a fixed sampling period $T$. The state is propagated according to
\[
x_{k+1} = x_k + Tf(x_k,u_k)
\]
The simulation is run over a finite time horizon, with the input voltage
periodically toggled to excite the system dynamics. At each time step,
the current state is stored to enable analysis and visualization of
the system behavior over time.

\section*{Results}
The time histories of the inductor current and capacitor voltage are shown
in the figure below. The system exhibits oscillatory behavior due to energy
exchange between the inductor and capacitor, with damping introduced by 
the resistor. \\
\begin{figure}[h]
    \centering
    \includegraphics[width=\linewidth]{../figures/states.pdf}
    \caption{Inductor current and capacitor voltage over time}
\end{figure}

\pagebreak

Next, the resistor and capacitor currents are consistent with the underlying
circuit relationships. \\
\begin{figure}[h]
    \centering
    \includegraphics[width=\linewidth]{../figures/current.pdf}
    \caption{Resistor and capacitor currents over time}
\end{figure}

Lastly, the equivalent impedance of the circuit. The variation reflects the 
time-dependent relationship between voltage and current in the system. \\
\begin{figure}[h]
    \centering
    \includegraphics[width=\linewidth]{../figures/impedance.pdf}
    \caption{Equivalent impedance over time}
\end{figure}

\section*{Conclusion}
This project demonstrates a simple framework for simulating continuous-time
dynamical systems using a state-space representation and forward Euler
integration. The same structure can be extended to more complex systems and
forms the basis for later state estimatiom methods.

\end{document}